\DocumentMetadata{
  pdfversion=2.0,
  pdfstandard=ua-2,
  lang=en-US,
  tagging=on,
  tagging-setup={math/setup=mathml-SE,tabsorder=structure}
}
\documentclass[11pt,letterpaper]{article}
\usepackage[margin=0.68in,includefoot]{geometry}
\usepackage{newcomputermodern}
\usepackage{amsmath,amscd,mathtools}
\usepackage{tikz,adjustbox}
\providecommand{\square}{\mdlgwhtsquare}
\usepackage[hidelinks]{hyperref}
\hypersetup{
  pdftitle={Topology First-Year Exam, January 2023, Part 1},
  pdfauthor={University of Florida Department of Mathematics},
  pdfsubject={Graduate mathematics examination},
  pdfdisplaydoctitle=true,
  bookmarksopen=true
}
\setcounter{secnumdepth}{0}
\setlength{\parindent}{0pt}
\setlength{\parskip}{0.28em}
\setlength{\emergencystretch}{3em}
\setlength{\leftmargini}{2.15em}
\setlength{\leftmarginii}{2.65em}
\setlength{\leftmarginiii}{3em}
\renewcommand{\labelenumi}{\textbf{\theenumi.}}
\pagestyle{empty}
\raggedbottom
\allowdisplaybreaks
\newenvironment{examproblems}[1]{%
  \begin{enumerate}%
  \setcounter{enumi}{#1}%
  \setlength{\topsep}{0.35em}%
  \setlength{\partopsep}{0pt}%
  \setlength{\itemsep}{0.48em}%
  \setlength{\parsep}{0.18em}%
}{\end{enumerate}}
\newenvironment{examsubparts}{%
  \begin{enumerate}%
  \setlength{\topsep}{0.18em}%
  \setlength{\partopsep}{0pt}%
  \setlength{\itemsep}{0.16em}%
  \setlength{\parsep}{0.1em}%
}{\end{enumerate}}
\newenvironment{examinstructions}{%
  \par\begingroup\itshape
}{%
  \par\endgroup\vspace{0.18em}
}
\makeatletter
\renewcommand\section{\@startsection{section}{1}{0pt}%
  {-1.25ex plus -.25ex minus -.1ex}%
  {0.55ex plus .1ex}%
  {\normalfont\large\bfseries}}
\renewcommand{\@maketitle}{%
  \newpage\null\vskip -1em%
  {\footnotesize\bfseries
    \href{https://www.ufl.edu/}{University of Florida}, \href{https://math.ufl.edu/}{Department of Mathematics}\par}%
  \vskip -0.5em%
  \tagstructbegin{tag=Title}%
    {\LARGE\bfseries\href{https://gma.math.ufl.edu/exams/topology-first-year/}{Topology First-Year Exam}, January 2023, Part 1\par}%
  \tagstructend%
  \vskip 0.25em%
  \tagmcbegin{artifact}%
    \hrule height 1pt%
  \tagmcend%
  \vskip 1em%
}
\makeatother
\title{Topology First-Year Exam, January 2023, Part 1}
\author{}
\date{}
\begin{document}
\maketitle
\thispagestyle{empty}
\begin{examinstructions}
Work the following problems and show all work. Support all statements to the best of your ability.
\end{examinstructions}
\begin{examproblems}{0}
\item
Show that a connected metric space having more than one point is uncountable.
\item
Give definition of retraction. Show that a retraction is a quotient map.
\item
Is the set of rationals~\(\mathbb{Q}\) in the subspace topology locally compact?
\item
Show that the set of irrational numbers in~\(\mathbb{R}\) is a Baire space.
\item
Let \(f:X\to X\) be a map of a compact metric space \((X,d)\) to itself that satisfies the following condition \(d(f(x),f(y))<d(x,y)\) for all \(x\ne y\). Prove that~\(f\) is continuous. Show that~\(f\) has a unique fixed point.
\end{examproblems}
\begin{examinstructions}
Answer the following with complete definitions or statements or short proofs.
\end{examinstructions}
\begin{examproblems}{5}
\item
State the Intermediate Value Theorem.
\item
Show that any continuous map \(f:S^2\to(0,1)\) of the 2-sphere \(S^2\) to the interval \((0,1)\) is neither surjective nor injective.
\item
State the Cantor–Schröder–Bernstein Theorem.
\item
Is the space \(\mathbb{R}_\ell\) connected? path connected?
\item
What is a basis of a topology? Does the set \(\{(a,b]\cup[c,d)\mid a<b<c<d\}\) form a basis of a topology on~\(\mathbb{R}\)?
\end{examproblems}
\end{document}
